\section{Data}
\label{sec:data}

All series are retrieved programmatically, checksummed, and rebuilt from raw
inputs. No value in this paper is typed by hand: every number in the text is a
macro generated from the analysis dataset, and every table and figure is written
by the pipeline that produced it.

\subsection{Statutory minimum wage}

The statutory series is parsed from the official history published by the
Direção-Geral do Emprego e das Relações de Trabalho (DGERT), which lists one row
per legal act with its effective date and the value it set. The series begins
with Decreto-Lei n.º 217/74 of 27 May 1974, which introduced a minimum wage of
3\,300 escudos, and runs to the present.

Portugal did not have a single minimum wage until 2005. Up to 1991 three
legally distinct wages coexisted --- domestic service, agriculture, and a general
regime covering the remaining activities --- falling to two between 1992 and 2004.
The general regime is the series conventionally quoted as the Portuguese minimum
wage, and it is the one used throughout. The special regimes were always set at
or below it, which the pipeline asserts on every date where regimes coexist.

Escudo amounts are converted at 200.482 escudos per euro, the irrevocable rate
fixed for the changeover. The act effective in January 2002 states both
currencies, and its published euro figure is preferred over the converted one.

Two independent checks validate the parsed series. First, the provider states
each act's percentage increase over the wage it replaced; recomputing that from
consecutive parsed values reconciles amounts, ordering, and regime assignment
simultaneously. Second, Eurostat compiles the same statutory wage from its own
reading of national law. Eurostat expresses it on a twelve-month basis whereas
Portuguese law states a monthly wage paid fourteen times a year, so its figure
equals the monthly level times $14/12$; undoing that convention reproduces the
parsed series to within a euro across the years checked.

The reconciliation exposes one gap in the national source: the act effective in
2000 is absent from the published history, while the 2001 entries still state
their increase relative to it. Because a year with no listed act is otherwise
indistinguishable from a year in which the wage was genuinely frozen --- as it was
in 2012 and 2013 --- the pipeline resolves the ambiguity against Eurostat, and
uses the Eurostat level only where the two compilers materially disagree. The
2000 observation is therefore attributed to Eurostat rather than to the national
legal source, and is flagged as such in the dataset.

\subsection{Prices and productivity}

Consumer prices are the World Bank index for Portugal, available from 1960.
Eurostat's harmonised index begins in 1996, two decades after the minimum wage
was introduced, and would truncate the historical layer at the point where it is
least informative.

Labour productivity is real gross domestic product per person employed from the
European Commission AMECO database (series \texttt{PRT.1.1.0.0.RVGDE}),
available from 1960. The Eurostat national-accounts equivalent begins in 1995
for Portugal. Over the 31 years in which the two overlap they are effectively
the same series: their annual growth rates correlate at unity and differ by at
most 0.003 percentage points. AMECO carries Commission projections beyond the
last outturn, which are excluded; only observed years enter the dataset.

\subsection{Coverage and exposure}

The exposure design requires the pre-existing share of workers paid at the
minimum wage, by region and industry. The Gabinete de Estratégia e Planeamento
(GEP) published these statistics, but has since withdrawn the host: every path
under its former domain now answers with an unrelated landing page. The
publications are retrieved instead from the Internet Archive and are registered
with their archival addresses and checksums, so the citation remains verifiable.

The recovered monitoring reports give coverage by economic activity. They do not
give a systematic breakdown by NUTS II region, which the identification strategy
in Section~\ref{sec:methods} requires. Constructing a region-by-industry measure
from these sources would require assuming that within-industry coverage does not
vary across regions, so that all regional variation comes from industry mix.
That assumption is not innocuous in a country whose regions differ sharply in
tourism intensity, and it is not made here. The price pass-through estimates are
therefore not reported in this version.
